\section{No Directive Moves the Cut}
\label{sec:ch13-no-directive-moves-the-cut}

\index{federation directives!page ownership|(}%
The next place to look is the directive vocabulary of
chapter~\ref{ch:federation-model}, because every one of those directives
describes what a subgraph can reach or depends on, and this looks like a
dependency. I went through the four this graph could use, and none of them
hands Sessions a score before its list resolver
runs~\autocite{apollodirectives}.

\index{requires@\code{"@requires}!does not reach a root field}%
\code{@requires} is the tempting one, because
chapter~\ref{ch:second-third-subgraph} used it for something that looks close:
\code{feedbackUrl} in Ratings is built from a title Sessions owns, and the
plan collected the title into the representation before the field resolver
ran. That is a field resolver being handed a value about the object it is
already resolving. It does nothing for a root field that has to choose its
objects first, and there is no representation of a session to attach a score
to until Sessions has chosen the session.

\index{provides@\code{"@provides}!needs the value already present}%
\code{@provides} points the other way. It lets a resolver return a field it
does not own, at one path, on the condition that the service already holds
the value; the directive neither fetches it nor keeps it current, and holding
the value means keeping a copy of a column in a service that does not own it,
which is why chapter~\ref{ch:second-third-subgraph} shipped none.
\code{@shareable} says two subgraphs can both resolve a field, and
copies nothing into either of them. If Sessions genuinely maintained an
average score of its own, marking both definitions shareable would describe
that, and it would be the local column and not the annotation that made
filtering possible. And \code{@listSize} is the cost analyzer's description of
how big a list may get, which chapter~\ref{ch:composition} already met and
switched off; it says nothing about how the list is ordered.

\index{search!is possible, and needs an owner}%
Those are claims about this schema and not a verdict on federation. What a
search surface needs is an owner that holds every value it filters and sorts
by at the moment it cuts the page, and this graph does not have one yet.

\index{search!where to put the owner}%
\index{projection!the other sense of the word}%
There are two places it could go, and I built both. A complete projection
inside Ratings, and projection here means a copy of the values to search by
rather than chapter~\ref{ch:data-without-n-plus-1}'s narrowing of columns,
produces the same page, the same shape of SQL, the same statement counts and
the same router plan as a separate service, and saves a process. It also
makes the service that stores scores the owner of a second domain, discovery,
that has nothing to do with writing a score. The other place is a fourth
subgraph, Search, which owns no session and no rating and holds a document
built from both for one purpose, finding sessions. That costs a process, and
it costs whatever keeps the document fresh, and I take it, because hiding
those two bills inside Ratings would not have made them smaller.
\index{federation directives!page ownership|)}%
